% GENERATED FILE. Edit canonical lessons, then run npm run generate:books.
% canonical-source-hash: fnv1a64:f923737a75cc12dc
% canonical-lessons: TA-C376-poratu, TA-C376-arppattam, TA-C376-perani, TA-C376-koru, TA-C376-manu

\chapter{Making Demands}
\label{ch:ta-making-demands}
\hlchaptermodality{376}

\begin{chapteropening}
\textbf{By the end of this chapter:} \emph{I can talk about protests, rallies and petitions.}

Complete the last lesson of chapter 376: I can talk about protests, rallies and petitions.
\end{chapteropening}

\section[\texorpdfstring{\ta{போராடு} (pōrāṭu)}{போராடு (pōrāṭu)}]{\ta{போராடு} (pōrāṭu) — to struggle; to fight for}

\label{lesson:TA-C376-poratu}



\begin{quote}
Before the new one: say the Tamil for influence, then the Tamil for to attract.
\end{quote}

\subsection*{You'll want to know: \ta{போராடு}}
\textbf{\ta{போராடு}} — \emph{pōrāṭu} — \textquotedblleft{}to struggle; to fight for\textquotedblright{}.

\begin{quote}
\ta{விவசாயிகள்} \ta{தங்கள்} \ta{உரிமைகளுக்காகப்} \ta{போராடுகிறார்கள்}.
\end{quote}

\begin{quote}
\emph{vivacāyikaḷ taṅkaḷ urimaikaḷukkākap pōrāṭukiṟārkaḷ.}
\end{quote}

\begin{quote}
The farmers are fighting for their rights.
\end{quote}

Past (I …): \ta{போராடினேன்}, \emph{pōrāṭiṉēṉ}.

X-\ta{க்காகப்} \ta{போராடு}, to fight for X. \ta{போராட்டம்}, a struggle or a protest.

\begin{cousinweb}
\textbf{\ta{போராடு}} (\emph{pōrāṭu}) is native Dravidian: DEDR 4540 lists Tamil \emph{poru}, \textquotedblleft{}to fight\textquotedblright{}, with \emph{pōr-āṭu}, \textquotedblleft{}to fight\textquotedblright{}, and \emph{pōr-āṭṭam}, \textquotedblleft{}struggle\textquotedblright{}.
\end{cousinweb}

\begin{grammarlens}[title={What you've built}]
One more word for giving reasons and opinions.
\end{grammarlens}

\subsection*{Your turn}
\begin{itemize}
\raggedright
  \item \emph{Say it:} \emph{pōrāṭu}
  \item \emph{Say it:} \emph{pōrāṭu}, once more
  \item \emph{Say it:} \emph{vivacāyikaḷ taṅkaḷ urimaikaḷukkākap pōrāṭukiṟārkaḷ.} — The farmers are fighting for their rights.
  \item \emph{Recall:} say the Tamil for influence, then the Tamil for to attract, then say \emph{pōrāṭu} again
\end{itemize}

\begin{tcolorbox}[breakable,colback=teal!4,colframe=teal!35!black,title={Before you move on}]
What does \ta{போராடு} mean? (\textquotedblleft{}To struggle; to fight for\textquotedblright{}.) Say it once more.
\end{tcolorbox}
\section[\texorpdfstring{\ta{ஆர்ப்பாட்டம்} (ārppāṭṭam)}{ஆர்ப்பாட்டம் (ārppāṭṭam)}]{\ta{ஆர்ப்பாட்டம்} (ārppāṭṭam) — a demonstration; a protest}

\label{lesson:TA-C376-arppattam}



\begin{quote}
Before the new one: say the Tamil for to attract, then the Tamil for to struggle.
\end{quote}

\subsection*{You'll want to know: \ta{ஆர்ப்பாட்டம்}}
\textbf{\ta{ஆர்ப்பாட்டம்}} — \emph{ārppāṭṭam} — \textquotedblleft{}a demonstration; a protest\textquotedblright{}.

\begin{quote}
\ta{மாணவர்கள்} \ta{கல்லூரிக்கு} \ta{முன்னால்} \ta{ஆர்ப்பாட்டம்} \ta{செய்தார்கள்}.
\end{quote}

\begin{quote}
\emph{māṇavarkaḷ kallūrikku muṉṉāl ārppāṭṭam ceytārkaḷ.}
\end{quote}

\begin{quote}
The students held a demonstration in front of the college.
\end{quote}

\ta{ஆர்ப்பாட்டம்} \ta{செய்}, to demonstrate. It can also mean a noisy show.

\begin{cousinweb}
\textbf{\ta{ஆர்ப்பாட்டம்}} (\emph{ārppāṭṭam}) is native, built on \emph{ārppu}, \textquotedblleft{}loud, tumultuous noise\textquotedblright{}, from \emph{ār}, \textquotedblleft{}to shout, roar\textquotedblright{} (DEDR 367): a gathering that makes itself heard.
\end{cousinweb}

\begin{grammarlens}[title={What you've built}]
One more word for giving reasons and opinions.
\end{grammarlens}

\subsection*{Your turn}
\begin{itemize}
\raggedright
  \item \emph{Say it:} \emph{ārppāṭṭam}
  \item \emph{Say it:} \emph{ārppāṭṭam}, once more
  \item \emph{Say it:} \emph{māṇavarkaḷ kallūrikku muṉṉāl ārppāṭṭam ceytārkaḷ.} — The students held a demonstration in front of the college.
  \item \emph{Recall:} say the Tamil for to attract, then the Tamil for to struggle, then say \emph{ārppāṭṭam} again
\end{itemize}

\begin{tcolorbox}[breakable,colback=teal!4,colframe=teal!35!black,title={Before you move on}]
What does \ta{ஆர்ப்பாட்டம்} mean? (\textquotedblleft{}A demonstration; a protest\textquotedblright{}.) Say it once more.
\end{tcolorbox}
\section[\texorpdfstring{\ta{பேரணி} (pēraṇi)}{பேரணி (pēraṇi)}]{\ta{பேரணி} (pēraṇi) — a rally; a march}

\label{lesson:TA-C376-perani}



\begin{quote}
Before the new one: say the Tamil for to struggle, then the Tamil for a demonstration.
\end{quote}

\subsection*{You'll want to know: \ta{பேரணி}}
\textbf{\ta{பேரணி}} — \emph{pēraṇi} — \textquotedblleft{}a rally; a march\textquotedblright{}.

\begin{quote}
\ta{ஆயிரக்கணக்கான} \ta{மக்கள்} \ta{பேரணியில்} \ta{கலந்துகொண்டார்கள்}.
\end{quote}

\begin{quote}
\emph{āyirakkaṇakkāṉa makkaḷ pēraṇiyil kalantukoṇṭārkaḷ.}
\end{quote}

\begin{quote}
Thousands of people took part in the rally.
\end{quote}

\ta{பேர்}, great + \ta{அணி}, a row: a great line of people.

\begin{cousinweb}
\textbf{\ta{பேரணி}} (\emph{pēraṇi}) is native: \emph{pēr}, \textquotedblleft{}great\textquotedblright{} (DEDR 4411), and \emph{aṇi}, \textquotedblleft{}array, row\textquotedblright{} (DEDR 117): a great procession.
\end{cousinweb}

\begin{grammarlens}[title={What you've built}]
One more word for giving reasons and opinions.
\end{grammarlens}

\subsection*{Your turn}
\begin{itemize}
\raggedright
  \item \emph{Say it:} \emph{pēraṇi}
  \item \emph{Say it:} \emph{pēraṇi}, once more
  \item \emph{Say it:} \emph{āyirakkaṇakkāṉa makkaḷ pēraṇiyil kalantukoṇṭārkaḷ.} — Thousands of people took part in the rally.
  \item \emph{Recall:} say the Tamil for to struggle, then the Tamil for a demonstration, then say \emph{pēraṇi} again
\end{itemize}

\begin{tcolorbox}[breakable,colback=teal!4,colframe=teal!35!black,title={Before you move on}]
What does \ta{பேரணி} mean? (\textquotedblleft{}A rally; a march\textquotedblright{}.) Say it once more.
\end{tcolorbox}
\section[\texorpdfstring{\ta{கோரு} (kōru)}{கோரு (kōru)}]{\ta{கோரு} (kōru) — to demand; to call for}

\label{lesson:TA-C376-koru}



\begin{quote}
Before the new one: say the Tamil for a demonstration, then the Tamil for a rally.
\end{quote}

\subsection*{You'll want to know: \ta{கோரு}}
\textbf{\ta{கோரு}} — \emph{kōru} — \textquotedblleft{}to demand; to call for\textquotedblright{}.

\begin{quote}
\ta{தொழிலாளர்கள்} \ta{அதிக} \ta{சம்பளம்} \ta{கோரினார்கள்}.
\end{quote}

\begin{quote}
\emph{toḻilāḷarkaḷ atika campaḷam kōriṉārkaḷ.}
\end{quote}

\begin{quote}
The workers demanded higher pay.
\end{quote}

Past (I …): \ta{கோரினேன்}, \emph{kōriṉēṉ}.

The noun is \ta{கோரிக்கை}, a demand or a request.

\begin{cousinweb}
\textbf{\ta{கோரு}} (\emph{kōru}) is native Dravidian: DEDR 2232 lists Tamil \emph{kōru}, \textquotedblleft{}to request, wish\textquotedblright{}, and \emph{kōrikkai}, \textquotedblleft{}request\textquotedblright{}.
\end{cousinweb}

\begin{grammarlens}[title={What you've built}]
One more word for giving reasons and opinions.
\end{grammarlens}

\subsection*{Your turn}
\begin{itemize}
\raggedright
  \item \emph{Say it:} \emph{kōru}
  \item \emph{Say it:} \emph{kōru}, once more
  \item \emph{Say it:} \emph{toḻilāḷarkaḷ atika campaḷam kōriṉārkaḷ.} — The workers demanded higher pay.
  \item \emph{Recall:} say the Tamil for a demonstration, then the Tamil for a rally, then say \emph{kōru} again
\end{itemize}

\begin{tcolorbox}[breakable,colback=teal!4,colframe=teal!35!black,title={Before you move on}]
What does \ta{கோரு} mean? (\textquotedblleft{}To demand; to call for\textquotedblright{}.) Say it once more.
\end{tcolorbox}
\section[\texorpdfstring{\ta{மனு} (maṉu)}{மனு (maṉu)}]{\ta{மனு} (maṉu) — a petition}

\label{lesson:TA-C376-manu}



\begin{quote}
Before the new one: say the Tamil for a rally, then the Tamil for to demand.
\end{quote}

\subsection*{You'll want to know: \ta{மனு}}
\textbf{\ta{மனு}} — \emph{maṉu} — \textquotedblleft{}a petition\textquotedblright{}.

\begin{quote}
\ta{கிராம} \ta{மக்கள்} \ta{ஆட்சியரிடம்} \ta{ஒரு} \ta{மனு} \ta{கொடுத்தார்கள்}.
\end{quote}

\begin{quote}
\emph{kirāma makkaḷ āṭciyariṭam oru maṉu koṭuttārkaḷ.}
\end{quote}

\begin{quote}
The villagers handed a petition to the district collector.
\end{quote}

\ta{மனு} \ta{கொடு}, to hand in a petition. Plural \ta{மனுக்கள்}.

\begin{cousinweb}
\textbf{\ta{மனு}} (\emph{maṉu}), \textquotedblleft{}a petition\textquotedblright{}, is an everyday word of Tamil offices; its origin is not settled.
\end{cousinweb}

\begin{grammarlens}[title={What you've built}]
One more word for giving reasons and opinions.
\end{grammarlens}

\subsection*{Your turn}
\begin{itemize}
\raggedright
  \item \emph{Say it:} \emph{maṉu}
  \item \emph{Say it:} \emph{maṉu}, once more
  \item \emph{Say it:} \emph{kirāma makkaḷ āṭciyariṭam oru maṉu koṭuttārkaḷ.} — The villagers handed a petition to the district collector.
  \item \emph{Recall:} say the Tamil for a rally, then the Tamil for to demand, then say \emph{maṉu} again
\end{itemize}

\begin{tcolorbox}[breakable,colback=teal!4,colframe=teal!35!black,title={Before you move on}]
What does \ta{மனு} mean? (\textquotedblleft{}A petition\textquotedblright{}.) Say it once more.
\end{tcolorbox}
